%% Section 6: Generalization Evidence Across Participants, Sessions, and Languages

\section{Generalization Evidence}
\label{sec:generalization}

\subsection{The Generalization Gap: A Quantitative Overview}

Of the 71 included studies, 40 (56.3\%) demonstrate \textit{no generalization} beyond the participant and session on which the model was trained. An additional 18 studies (25.4\%) conduct multi-session evaluations within a single participant (typically 1-3 additional sessions in sequence, without explicit long-term follow-up). Only 13 studies (18.3\%) demonstrate cross-participant generalization \cite{Peterson2021,deBorman2024}. Cross-language generalization appears in exactly 2 studies (2.8\%). Long-term cross-session stability with performance retention quantified across weeks or months appears in exactly 1 study (1.4\%) \cite{Ye2025}.

This generalization gap is the field's highest-priority bottleneck for clinical translation. Systems that work beautifully in a single participant's session fail immediately with new participants or new sessions. For a technology whose clinical promise depends on deployment in clinics across diverse patients and daily use, this is a critical limitation.

Cross-participant generalization remains the field's highest-priority bottleneck for clinical translation. Of the 71 included studies, only 13 (18.3\%) demonstrate that a model trained on one participant maintains reasonable performance on a different participant. The majority of cross-participant studies (n=13) employ transfer learning, manifold alignment, or domain adaptation techniques borrowed from computer vision literature but adapted for neural signal processing.

Transfer learning approaches leverage the observation that early layers of neural networks learn modality-generic features (e.g., spectro-temporal patterns in EEG), while later layers specialize to individual differences. Peterson et al. \cite{Peterson2021} demonstrated this principle by training a shared neural representation across 28 EEG participants: features extracted from the first layers generalized across participants, while final output layers required participant-specific adaptation. Cross-participant performance retained 60-80\% of single-participant accuracy, with 80\% retention on simple tasks (vowel classification, 4-class) and 60\% retention on harder tasks (continuous speech, 125K-vocabulary). The key insight: individual neural motor encoding differs substantially, but the underlying neural-to-acoustic mapping contains participant-invariant structure exploitable via multi-task learning.

Manifold alignment offers a complementary approach. Neural activity in single participants traces a low-dimensional manifold (typically 5-50 dimensions) embedded in high-dimensional recording space (64-1024 dimensions for ECoG, 64-256 for EEG). Different participants' manifolds have different orientations and curvatures in neural space, even though the underlying computation (speech motor planning) should be identical. Peterson et al. \cite{Peterson2021} applied manifold alignment logic across participants, achieving 60-80\% transfer. Domain adaptation methods (Singh et al. \cite{Singh2025}, deBorman et al. \cite{deBorman2024}) bridge this further by learning invariant representations that eliminate participant-specific signal characteristics (electrode noise, impedance, individual alpha frequency). Singh et al.'s distributed sEEG network across 25 participants at multiple clinical sites achieved 50-75\% of single-site performance, suggesting that domain-invariant features are extractable even across clinical site variation. Methodologically, manifold alignment approaches have been validated in non-human primate models, but human speech system applicability remains limited to the few cross-participant generalization studies noted above.

Performance retention rates cluster into two regimes. Simpler tasks (binary classification, vowel identification, event detection) show 70-90\% cross-participant transfer. Complex tasks (open-vocabulary continuous speech, real-time communication) show 50-70\% transfer. The gap reflects the fact that complex tasks encode more participant-specific information: motor planning for unrestricted speech leverages individual vocabulary, accent, speaking rate, idiolect; decoding systems must learn these individual patterns for optimal performance. Transfer learning sacrifices within-participant performance to achieve cross-participant generalization - a fundamental tradeoff.

Critical constraints limit practical deployment of cross-participant methods. Most require large multi-participant datasets ($N = 8$-$28$ participants) for training, yet the clinical patient population is scattered across dozens of hospitals. Aggregating 28 ALS participants across sites demands coordinated data collection protocols and privacy infrastructure currently absent in the field. Furthermore, transfer learning requires careful cross-validation: naive train-test splits that leak information across participants inflate generalization estimates by 5-15 percentage points \cite{Peterson2021}. Only 4 of 13 cross-participant studies employed explicit participant-level cross-validation; others used within-participant cross-validation and claimed cross-participant generalization - a common pitfall inflating perceived generalization.

\subsection{Cross-Participant Generalization (n=13)}

Three landmark studies exemplify successful cross-participant transfer, each revealing different technical pathways and constraints.

Peterson et al. \cite{Peterson2021} achieved cross-participant generalization across 28 EEG participants using neural manifold alignment. Participants performed a constrained speech production task (imagined and overt speech from a limited vocabulary); researchers extracted neural activity from sensorimotor regions. The approach: (1) Learn a shared manifold embedding across all 28 participants using autoencoders; (2) Learn participant-specific transformation matrices mapping each participant's neural space to the shared manifold; (3) Train a decoder on the shared manifold. Cross-participant evaluation (leave-one-participant-out cross-validation) achieved 60-80\% of within-participant performance on vowel classification, 50-70\% on consonant classification, and 50\% on continuous word decoding. The key success factor: EEG signals contain substantial shared structure across participants (common frequency bands, common gross topography), making manifold alignment tractable. The limitation: 28 participants required; scaling to hundreds remains computationally expensive.

Singh et al. \cite{Singh2025} demonstrated cross-modal, multi-site generalization using distributed sEEG recordings across 25 participants implanted at three different clinical sites (UCSF, Stanford, MGH). Sites used different electrode manufacturers, different reference schemes, and different amplifiers - introducing systematic site-specific variation (``batch effects''). Approach: (1) Standardize signal preprocessing across sites via harmonization techniques borrowed from multi-site neuroimaging; (2) Train a group decoder on pooled data from all sites; (3) Evaluate generalization to held-out participants at each site. Results: Group decoder trained on 18 participants from two sites achieved 50-75\% of site-specific single-participant performance when tested on held-out participants from the third site. Clinical relevance: This demonstrates that standardized protocols can enable multi-site deployment, critical for moving from single-lab demonstrations to clinical networks.

deBorman et al. \cite{deBorman2024} showed that simpler tasks enable better cross-participant transfer. The task: Detect motor intention onset (did the participant attempt to speak?) from ECoG recordings across 8 participants. This binary classification task is substantially simpler than continuous speech decoding. Cross-participant performance: 70\% accuracy vs. baseline (chance = 50\%), retaining 85-90\% of within-participant accuracy. The success reflects two factors: (1) Motor intention is neurally consistent across participants (all motor cortices implement similar intention-detection circuits); (2) Binary classification has lower sample complexity than speech decoding, reducing data requirements. The limitation: Event detection alone is insufficient for clinical speech prosthesis; actual speech decoding remains harder to generalize.

Collectively, these three studies establish that cross-participant transfer is achievable but demands either (1) large datasets ($N = 8$-$28$), (2) careful signal harmonization, or (3) simplified tasks. Near-term clinical deployment will likely require combinations of all three: smaller patient cohorts will generalize only partially, requiring frequent participant-specific calibration; multi-site networks will standardize protocols to minimize batch effects; and early deployed systems will enable simpler tasks (yes/no responses, fixed-vocabulary communication) before tackling open-vocabulary speech.

\subsection{Cross-Session Generalization (Long-Term Stability: n=1)}

Only one human-participant study explicitly evaluated whether a model trained on one session maintains performance on subsequent sessions weeks or months later. Ye et al. \cite{Ye2025} demonstrated $>$80\% performance retention across 4 weeks of daily use with recalibration. As a methodological reference (though not part of the included 71 studies per PICOS criteria for human participants), Degenhart et al.'s work on non-human primates showed that manifold alignment can preserve 95\% performance across 6 months, suggesting that the technical approach is viable. However, applicability to human speech systems remains unvalidated. This extreme scarcity of human evidence (only 1 study out of 71) highlights a critical gap in generalization research.

Neural signals drift substantially over time due to multiple mechanisms. Electrode micromotion - gradual displacement of recording electrodes due to brain pulsation, inflammation, or mechanical stress - causes gradual signal attenuation \cite{Vansteensel2023}. Glial scarring around electrodes develops over weeks to months post-implantation (weeks 1-3 proliferation phase, week 6 stabilization with 200 \textmu{}m capsule formation) and increases electrode impedance and reduces signal amplitude \cite{GlialScar2014,Vansteensel2023}. Neural plasticity - reorganization of motor cortex representations following electrode implantation - can shift the tuning of recorded neurons over weeks. These biological processes combine to degrade decoder performance by 10-30\% per month without recalibration, depending on electrode type and implant duration \cite{Degenhart2020}.

Recalibration mechanisms address signal drift through several approaches. Online adaptation (Degenhart et al. approach) continuously updates decoder weights to track signal changes, maintaining 95\% performance across 6 months without explicit user training. Periodic recalibration (daily or weekly) retrains the decoder on fresh data from the current session, resetting performance to single-session levels at the cost of required calibration time. Adaptive learning uses Bayesian or kernel adaptation to smooth transitions between sessions while leveraging prior knowledge from earlier sessions. The choice between approaches depends on clinical constraints: daily recalibration is operationally feasible for home-based systems but burdensome; online adaptation requires stable adaptation algorithms that don't catastrophically forget.

Time scale evidence is limited but critical for clinical planning. Degenhart et al. demonstrated $>$95\% stability for 6 months with manifold alignment. The second cross-session study showed $>$80\% retention for 4 weeks with recalibration. No study has demonstrated $>$6 months stability without some form of recalibration or adaptation. Only 2 studies achieve $>$1 month long-term stability (distinct from the 18 studies conducting multi-session evaluation within days). This extreme scarcity means clinical deployment cannot rely on published evidence for anything beyond $\sim$2-3 months of continuous use. This constraint directly impacts clinical business models: patients with rapidly progressive ALS (life expectancy 2-5 years) can tolerate bimonthly recalibration, but longer-term stability is essential for stroke patients or other static paralysis conditions where lifelong use is expected.

Clinical implications of cross-session limitations are substantial. A deployed speech neuroprosthesis cannot simply be trained once and used indefinitely; it requires either (1) frequent user-supervised recalibration (daily, weekly, or monthly), or (2) sophisticated online adaptation algorithms that are still experimental. Daily recalibration is operationally feasible (analogous to calibrating a cochlear implant battery), but weekly or monthly recalibration demands high patient adherence and clinical follow-up. For ALS patients with limited remaining survival, bimonthly recalibration is acceptable; for younger patients or non-progressive conditions, this is a significant limitation. The two published long-term stability studies did not evaluate clinical usability or patient satisfaction with recalibration burden, representing another critical gap in deployment readiness.

Future research must prioritize cross-session generalization. Beyond the two existing studies, the field lacks evidence on: (1) What signal drift mechanisms dominate over time (electrode micromotion vs. glial scarring vs. neural plasticity)? (2) Can online adaptation algorithms maintain $>$90\% performance for 12+ months? (3) What is the minimal recalibration frequency for clinical acceptability? (4) Do cross-participant transfer learning and cross-session transfer learning interact (does pre-training on other participants improve cross-session stability)? Answering these questions is prerequisite for clinical translation.

\subsection{Cross-Language Generalization (n=2)}

Only two studies evaluated whether speech decoding models trained on one language transfer to other languages, representing 2.8\% of the included 71 studies and highlighting an extreme gap in generalization evidence. The first study, d'Ascoli et al. \cite{dAscoli2025}, trained a CNN-Transformer architecture on a massive multilingual dataset spanning 723 participants across three languages: English, French, and Dutch. The model achieved $>$40\% word-level accuracy (above chance for the respective language vocabularies) when tested on native speakers of each language, demonstrating partial cross-language transfer despite substantial phonological and phonotactic differences between the languages. The key success factor was dataset scale: 723 participants enabled the model to learn language-invariant motor-phonetic mappings that transferred across languages. This study suggests that motor planning contains participant-invariant and language-invariant structure exploitable through large-scale transfer learning.

The second cross-language study, Li et al. \cite{Li2022}, examined a qualitatively different transfer scenario using high-density surface electromyography (HD-sEMG) recordings rather than neural recordings. HD-sEMG measures motor muscle activation during speech preparation and execution, capturing the motor planning rather than cortical representation of speech. Li et al. trained a motor-speech encoder-decoder on Mandarin Chinese speech production (phonetically and tonally complex) and tested it on English (simpler phonotactic structure). The model retained 80\% of single-language performance when transferring from Mandarin to English, suggesting that phoneme-motor mapping is substantially language-invariant at the muscle level. This result contrasts with d'Ascoli findings (40\% transfer on cortical EEG) and suggests that motor-level representations may be more language-universal than cortical representations.

Why is cross-language generalization so rare despite potential clinical importance? First, data requirements are substantial: d'Ascoli required 723 multilingual participants to achieve 40\% transfer; most brain-to-speech labs operate with $<$30 participants total. Aggregating multilingual participants across sites requires international coordination and language-specific preprocessing expertise currently rare in the field. Second, language-specific preprocessing challenges are non-trivial: different languages have different phoneme inventories (English has 44 phonemes; Mandarin has 21 consonants, 5 vowels, and 4 tones forming 22 total phoneme units with tone) \cite{Mandarin2020,Li2022}, different neural organizational principles (tonal languages recruit pitch-processing areas; non-tonal languages don't), and different phonotactic constraints affecting neural encoding. Preprocessing pipelines optimized for English typically fail on other languages without retuning. Third, the clinical value proposition is unclear: most deployed systems target single-language populations (English-only in USA, French-only in Canada). Multilingual capability is valuable for immigrant communities or multilingual users but represents a smaller market than monolingual systems.

Cross-language transfer has important implications for motor planning universality. If phoneme-motor mappings are universal across languages (as Li et al. suggest), then portable HD-sEMG systems could theoretically support multilingual users with minimal retraining. Conversely, if cortical language organization is language-specific (as d'Ascoli's partial transfer suggests), then cortical systems would require language-specific calibration. The limited evidence suggests a hybrid picture: motor-level representations are partially language-universal (80\% transfer on sEMG); cortical representations are partially language-universal (40\% transfer on EEG). Full language-agnostic systems remain aspirational; practical deployment will likely require language-specific or language-aware models.

Future cross-language research should prioritize: (1) Does multilingual pre-training on cortical recordings (EEG/MEG/invasive) improve single-language performance through regularization? (2) Which phonetic features transfer across languages (manner of articulation, place of articulation, voicing) and which are language-specific? (3) Can phoneme-level transfer learning outperform word-level transfer learning for language generalization? (4) What is the minimal multilingual dataset size for $>$60\% transfer performance? Answering these questions would unlock assistive technology for multilingual populations currently underserved by monolingual systems.

\subsection{Cross-Modal Generalization (n=5)}

Five studies attempted to train decoding models on one recording modality and test on a different modality, seeking to leverage advantages of multiple modalities or enable transition from research (invasive) to clinical (non-invasive) platforms \cite{Peterson2021,deBorman2024}. Cross-modal generalization is appealing because it could theoretically allow systems trained on high-signal invasive recordings to transfer to portable non-invasive systems, dramatically improving clinical deployability. However, the persistent rarity of successful cross-modal transfer (only 5/71 studies) reflects fundamental technical and biological obstacles that have not been overcome.

The core technical challenge is modality mismatch in the feature space. Invasive intracortical arrays achieve spatial resolution of 10 micrometers with signal amplitudes of 100-1000 \textmu{}V; scalp EEG achieves 40-60 mm resolution with signal amplitudes of 10-100 \textmu{}V. This 1000-fold spatial resolution difference and 10-fold amplitude difference means that intracortical and EEG data occupy radically different feature spaces. Intracortical signals encode individual neuron tuning properties; EEG encodes aggregate population-level activity across large cortical regions. Training a decoder on intracortical data and testing on EEG is fundamentally asking the decoder to extract fine-grained neural tuning information from aggregated population signals - a lossy transformation with no algorithmic remedy. Domain adaptation approaches (adversarial domain alignment, CycleGAN-style feature transformation) can reduce but not eliminate this mismatch.

Performance retention in the few published cross-modal studies clusters around 50-70\%. The most successful cross-modal transfer within published literature (ECoG to EEG) achieved 50-65\% performance retention; this success likely reflects the modalities' greater similarity compared to intracortical-to-EEG pairs (both ECoG and EEG measure lower-frequency, higher-amplitude signals compared to intracortical spikes). Conversely, attempts to transfer from intracortical microelectrode arrays to EEG have achieved 30-40\% retention, suggesting the spatial resolution gap is fundamentally limiting. No published study has demonstrated successful intracortical-to-fMRI transfer (the combination would be particularly challenging due to temporal resolution incompatibility in addition to spatial mismatch).

Domain adaptation strategies employed in the five cross-modal studies include: (1) Feature alignment using adversarial training to learn modality-invariant representations \cite{Peterson2021,Ma2023,Zhang2026Align}; (2) Intermediate modality bridging (e.g., ECoG as intermediate between intracortical and EEG); (3) Multi-task learning to identify shared structure across modalities \cite{Yang2026Corteg}; (4) Transfer learning with fine-tuning on small cross-modal datasets \cite{deBorman2024,Liu2026CrossBrain}. None achieved dramatic success ($>$80\% retention). The limited evidence suggests that truly modality-invariant features (universal across invasive/non-invasive or high-resolution/low-resolution recordings) may not exist; instead, each modality captures sufficiently distinct aspects of neural speech encoding that transfer is partial and lossy.

Clinical implications of cross-modal limitations are significant. A clinician cannot simply record invasive data in the lab, develop a high-performance decoder, and deploy it on non-invasive EEG in the clinic \cite{Yang2026Corteg,Liu2026CrossBrain}. Cross-modal transfer would be the ideal goal for brain-to-speech deployment (enabling lab development on high-signal invasive systems, clinic deployment on portable non-invasive systems), but it remains unachieved. Clinical deployment instead requires modality-specific development: systems intended for EEG deployment must be developed and validated on EEG data; systems intended for intracortical arrays must be developed on intracortical data. This constraint means that the low-SNR, low-resolution challenges of non-invasive modalities cannot be escaped through domain adaptation - they must be addressed through algorithmic innovation within each modality.

Why has cross-modal transfer proven so difficult despite decades of domain adaptation research in computer vision and speech recognition? The answer lies in the biological and physical differences between modalities. In computer vision, different image modalities (RGB, infrared, depth) capture overlapping information at different wavelengths but retain spatial coherence and similar feature hierarchies. In brain-to-speech, different neural recording modalities measure fundamentally different quantities (action potentials, local field potentials, aggregate hemodynamic responses) at vastly different spatial-temporal resolutions. The feature distributions are sufficiently divergent that standard domain adaptation approaches reach saturation. Future cross-modal research should focus on: (1) Identifying modality-invariant features (e.g., directional tuning properties preserved across modalities); (2) Intermediate transfer (using high-resolution non-invasive modalities like fMRI as bridge to EEG); (3) Complementary multi-modal fusion (combining EEG and fMRI simultaneously) rather than cross-modal transfer. Direct invasive-to-non-invasive transfer is unlikely to succeed at clinically useful performance levels.

\subsection{Implications for Clinical Deployment}

The generalization evidence directly constrains deployment scenarios. Current evidence supports:
\begin{enumerate}
    \item \textbf{Personalized, single-session systems:} Feasible. Train a custom model for each patient in their first session. Performance will be high but will degrade over months as neural signals drift.
    
    \item \textbf{Multi-session stable systems:} Limited evidence. Only 2 studies show >1 month stability; no study shows >6 months. Daily recalibration likely required.
    
    \item \textbf{Cross-participant transfer systems:} Marginal evidence. 13 studies show partial transfer (50-90\% of single-participant performance with pre-training). Practical applicability uncertain.
    
    \item \textbf{Language-agnostic systems:} Minimal evidence. Only 2 studies; no commercial product demonstrates this. English dominance (94\% of studies) reflects both research focus and practical limitation.
\end{enumerate}

